\section{The length--area argument}\label{sec:rigidity}

The trace statement now supplies the line estimates needed for a flux
lower bound.  Comparing that bound with the source and target area
bounds will force the normalized conformal map to have derivative of
unit modulus and corresponding disks to have equal radii.

\begin{proof}[Proof of Theorem~\ref{thm:main}]
If $\Omega=\Chat$, its conformal image is also the sphere, by
compactness, and its automorphisms are M\"obius transformations.
Assume henceforth that the complement is nonempty.  Apply separate
M\"obius changes in source and target to move a domain point and its
image to infinity.  A further target affine change gives
\begin{equation}\label{eq:rigidity-normalization}
 f(z)=z+\sum_{n\ge1}\alpha_nz^{-n}
 \quad\text{near infinity}.
\end{equation}
These changes preserve circle domains and conformal removability.
All complementary components are now bounded Euclidean disks or
points.  Their correspondence is supplied by Lemma~\ref{lem:ends}.

Index by $j$ precisely those corresponding pairs for which at least
one component is nondegenerate.  Write
\[
 K_j=\overline D(a_j,r_j),\qquad
 K'_j=\overline D(b_j,s_j),
\]
allowing a zero radius to mean a singleton.  The index set is finite
or countable, and
\begin{equation}\label{eq:radius-summability}
 \sum_jr_j^2<\infty,\qquad \sum_js_j^2<\infty,\qquad
 \sum_jr_js_j<\infty.
\end{equation}
Indeed the positive-radius disks on each side have disjoint interiors
in a bounded region, and the last assertion follows by
Cauchy--Schwarz.  All other complementary components are common
points: both the source and target component are singletons.
Let $\tau$ denote their correspondence.

Choose $R$ strictly larger than a radius containing the source
complement.  Define a vector field $F$ on $\overline D_R$ by using
$f$ in $\Omega$, the center $b_j$ on $K_j$, and $\tau(p)$ at every
common point $p$.  This is a bounded measurable field.  Boundedness
follows from Lemma~\ref{lem:ends}, or from the bounded side of the
image outer contour.  Measurability follows since there are
countably many indexed closed components, the common-point set is
Borel, and its singleton trace map is continuous relative to that
set.  No continuity of $F$ at an indexed component is asserted.

\paragraph{A slice estimate.}
Remove disjoint slightly enlarged concentric disks around the first
$N$ indexed source components, using a positive enlargement also when
$r_j=0$.  Denote the remaining round region by $A_N$.  The sizes of
the enlargements will tend to zero with $N$ fixed.

Consider a closed component segment $J=[x,y]$ of an axial slice of
$A_N$, directed toward increasing coordinate.  At domain points use
the projected value of $f$; at a complementary component use the
entire projection interval of its corresponding target component.
These fibers are compact intervals in a fixed bounded interval, and
their relation has closed graph by Lemma~\ref{lem:ends}.  Their union
over $J$ is connected.  To verify this directly, a missing value $t$
between two attained values would place each fiber wholly below or
wholly above $t$.  Upper semicontinuity makes these two classes open
in $J$, giving a separation of the segment.

Consequently all values between the selected endpoint values are
covered by the projected images of domain arcs, common-point traces,
and the full intervals of met indexed components.  By
Proposition~\ref{prop:trace-null}, the common traces have length zero
on almost every such line.  The countably many domain arcs have
total projected-image length at most the integral of $|f'|$ over
them.  Each met indexed component contributes at most $2s_j$.
Writing $F_{\rm ax}$ for the coordinate along the slice, we obtain
\begin{equation}\label{eq:slice-estimate}
 F_{\rm ax}(y)-F_{\rm ax}(x)
 \le \int_{J\cap\Omega}|f'|\,\dd s
       +\sum_{\substack{j>N\\K_j\cap J\ne\varnothing}}2s_j
\end{equation}
for almost every line in either direction.  In fact the same argument
bounds the absolute endpoint difference.  Endpoints on an enlarged
circle use the same field $F$ for both slicing directions.

\paragraph{Multiplicity of the tail terms.}
For fixed $N$ the enlargements can be chosen so that each remaining
positive-radius source component meets at most one enlarged disk.
For $N\ge2$, let $\delta>0$ be the minimum separation of the first
$N$ source components.  Make their enlargements small enough to
retain gaps greater than $\delta/2$.  A remaining disk of diameter
less than $\delta/2$ cannot meet two of them.  Only finitely many
remaining disks have larger diameter, and their positive distances
from the first $N$ components permit further shrinking of the
enlargements so that these finitely many disks meet none.  This
argument includes point components among the first $N$.  For $N=1$
there is nothing to prove, and for $N=0$ nothing is removed.

A line meets a remaining disk in an interval; deleting at most one
enlarged-disk interval leaves at most two retained pieces.  Thus its
cost in \eqref{eq:slice-estimate} is counted at most twice.  Integrating
transversely gives at most $8r_js_j$ per direction and at most
$16r_js_j$ for both directions.  Indexed source point components are
missed by almost every axial line, even when their target components
are disks.  All exceptional indexed point lines form a countable set.
The tail estimates are integrable by \eqref{eq:radius-summability}.

\paragraph{Flux.}
Integrating \eqref{eq:slice-estimate} over horizontal and vertical
lines gives the flux through the outer circle minus the flux through
the enlarged inner circles.  This is the elementary projection
formula for arclength on a circle: the transverse projection Jacobians
are $|n_X|$ for horizontal slices and $|n_Y|$ for vertical slices.
It applies to bounded measurable boundary values.  The normal on an inner circle is taken outward
from its removed disk.  The two domain integrals are bounded by
$2\int_{\Omega\cap D_R}|f'|\,\dd A$.

On the outer circle $F=f$, and
\eqref{eq:rigidity-normalization} gives exactly
\begin{equation}\label{eq:outer-flux}
 \int_{\partial D_R}F\cdot n\,\dd s
 =\operatorname{Re}\int_0^{2\pi}
     f(Re^{it})e^{-it}R\,\dd t
 =2\pi R^2.
\end{equation}
On the enlarged circle about $K_j$, subtract the constant target
center $b_j$, whose flux is zero.  As its radius decreases to $r_j$,
Lemma~\ref{lem:ends} puts all its assigned values within distance
$s_j+o(1)$ of $b_j$, uniformly along the circle.  Hence its inner
outward flux $I_j$ satisfies
\begin{equation}\label{eq:inner-flux}
 \limsup I_j\le2\pi r_js_j.
\end{equation}
When $r_j=0$ the circle length tends to zero; when $s_j=0$ the target
displacement tends to zero.  Thus both types of mismatch are covered.

Moving the inner fluxes to the right and letting the finitely many
enlargements shrink, while retaining the preceding multiplicity
property, gives
\[
 2\pi R^2
 \le 2\int_{\Omega\cap D_R}|f'|\,\dd A
    +2\pi\sum_{j\le N}r_js_j
    +16\sum_{j>N}r_js_j.
\]
Let $N$ exhaust the indexed components.  In the finite or empty case
the corresponding terminal sum is already zero.  We conclude that
\begin{equation}\label{eq:limiting-flux}
 2\pi R^2\le2\int_{\Omega\cap D_R}|f'|\,\dd A
                      +2\pi\sum_jr_js_j.
\end{equation}

\paragraph{Area and equality.}
The source domain portion and the interiors of its complementary
disks are disjoint subsets of $D_R$, so
\begin{equation}\label{eq:source-area}
 \int_{\Omega\cap D_R}1\,\dd A+\pi\sum_jr_j^2\le\pi R^2.
\end{equation}
The image of $\partial D_R$ is a Jordan contour.  Its bounded side
contains $f(\Omega\cap D_R)$ and all target complementary components:
the compact exterior source disk is mapped homeomorphically onto the
closed Jordan side containing infinity.  Indeed its interior is mapped
openly, its boundary is the image contour, and compactness makes its
image closed in that side; connectedness then identifies the entire
side.  These pieces are disjoint.  By the Laurent
area formula its enclosed area is
\[
 \pi R^2-\pi\sum_{n\ge1}n|\alpha_n|^2R^{-2n}\le\pi R^2.
\]
Conformal change of variables therefore gives
\begin{equation}\label{eq:target-area}
 \int_{\Omega\cap D_R}|f'|^2\,\dd A
       +\pi\sum_js_j^2\le\pi R^2.
\end{equation}
No area-zero assertion about the target dust is needed: any area in
its uncountably many singleton components is simply omitted from
the left side of this inequality.

Adding \eqref{eq:source-area} and \eqref{eq:target-area}, then using
\eqref{eq:limiting-flux}, yields
\begin{equation}\label{eq:nonnegative-rigidity}
 \boxed{\displaystyle
 \int_{\Omega\cap D_R}(|f'|-1)^2\,\dd A
       +\pi\sum_j(r_j-s_j)^2\le0.}
\end{equation}
Both terms are nonnegative.  Continuity of the integrand shows that
$|f'|=1$ throughout the nonempty open set $\Omega\cap D_R$.  On any
disk compactly contained in this set, the Open Mapping Theorem makes
$f'$ constant, so $f(z)=\lambda z+c$ there.  The meromorphic Identity
Theorem on the connected domain $\Omega$ gives this equality
everywhere in $\Omega$.  Finally
\eqref{eq:rigidity-normalization} forces $\lambda=1$ and $c=0$.
Undoing the source and target changes of coordinates proves that the
original equivalence is the restriction of a M\"obius transformation.
\end{proof}
